% Lines 4596-4967 from original attention_transformers_notes_ghost.tex
% Appendix: Practical Implementation Guide

%==============================================================================
\section{Practical Implementation Guide}
%==============================================================================

\subsection{Building Production-Ready Transformers}

\subsubsection{Efficient Key-Value Caching}

For autoregressive generation, caching previous key-value pairs is crucial:

\begin{lstlisting}[language=Python, caption=Key-Value Cache Implementation]
@dataclass
class KVCache:
    """Cache for storing key-value pairs across generation steps"""
    keys: list[Tensor]
    values: list[Tensor]
    max_seq_len: int
    current_seq_len: int = 0

class CachedMultiHeadAttention(nn.Module):
    """Multi-head attention with KV caching for generation"""

    def __init__(
        self,
        d_model: int,
        n_heads: int,
        max_seq_len: int,
        dropout: float = 0.1
    ) -> None:
        super().__init__()
        self.d_model = d_model
        self.n_heads = n_heads
        self.d_k = d_model // n_heads
        self.max_seq_len = max_seq_len

        # Projections
        self.w_q = nn.Linear(d_model, d_model)
        self.w_k = nn.Linear(d_model, d_model)
        self.w_v = nn.Linear(d_model, d_model)
        self.w_o = nn.Linear(d_model, d_model)

        self.dropout = nn.Dropout(dropout)
        self.scale = 1.0 / math.sqrt(self.d_k)

        # Initialize cache
        self.cache: KVCache | None = None

    def init_cache(self, batch_size: int) -> None:
        """Initialize KV cache for generation"""
        self.cache = KVCache(
            keys=[],
            values=[],
            max_seq_len=self.max_seq_len,
            current_seq_len=0
        )

        # Pre-allocate cache tensors for each layer
        for _ in range(self.n_heads):
            self.cache.keys.append(
                torch.zeros(
                    batch_size, self.max_seq_len, self.d_k,
                    device=next(self.parameters()).device
                )
            )
            self.cache.values.append(
                torch.zeros(
                    batch_size, self.max_seq_len, self.d_k,
                    device=next(self.parameters()).device
                )
            )

    def forward(
        self,
        query: Tensor,
        key: Tensor | None = None,
        value: Tensor | None = None,
        mask: Tensor | None = None,
        use_cache: bool = False
    ) -> tuple[Tensor, KVCache | None]:
        batch_size, seq_len, _ = query.shape

        # Project Q, K, V
        q = self.w_q(query).view(batch_size, seq_len, self.n_heads, self.d_k)
        q = q.transpose(1, 2)  # (batch, n_heads, seq_len, d_k)

        if key is None:
            key = query
        if value is None:
            value = query

        k = self.w_k(key).view(batch_size, -1, self.n_heads, self.d_k)
        v = self.w_v(value).view(batch_size, -1, self.n_heads, self.d_k)
        k = k.transpose(1, 2)
        v = v.transpose(1, 2)

        # Use cache if enabled
        if use_cache and self.cache is not None:
            # Update cache
            cache_seq_len = self.cache.current_seq_len

            # Store new keys and values
            for head_idx in range(self.n_heads):
                self.cache.keys[head_idx][:, cache_seq_len:cache_seq_len + seq_len] = k[:, head_idx]
                self.cache.values[head_idx][:, cache_seq_len:cache_seq_len + seq_len] = v[:, head_idx]

            self.cache.current_seq_len += seq_len

            # Use cached keys and values
            k_list = []
            v_list = []
            for head_idx in range(self.n_heads):
                k_list.append(
                    self.cache.keys[head_idx][:, :self.cache.current_seq_len]
                )
                v_list.append(
                    self.cache.values[head_idx][:, :self.cache.current_seq_len]
                )

            k = torch.stack(k_list, dim=1)
            v = torch.stack(v_list, dim=1)

        # Compute attention scores
        scores = torch.matmul(q, k.transpose(-2, -1)) * self.scale

        # Apply mask if provided
        if mask is not None:
            scores = scores.masked_fill(mask == 0, -1e9)

        # Softmax
        attn_weights = F.softmax(scores, dim=-1)
        attn_weights = self.dropout(attn_weights)

        # Apply attention to values
        context = torch.matmul(attn_weights, v)

        # Reshape and project output
        context = context.transpose(1, 2).contiguous()
        context = context.view(batch_size, seq_len, self.d_model)
        output = self.w_o(context)

        return output, self.cache
\end{lstlisting}

\subsubsection{Distributed Training Strategies}

Implement model parallelism for large transformers:

\begin{lstlisting}[language=Python, caption=Model Parallel Transformer]
class ModelParallelTransformerLayer(nn.Module):
    """Transformer layer with model parallelism"""

    def __init__(
        self,
        d_model: int,
        n_heads: int,
        d_ff: int,
        world_size: int,
        rank: int
    ) -> None:
        super().__init__()
        self.d_model = d_model
        self.n_heads = n_heads
        self.world_size = world_size
        self.rank = rank

        # Divide model dimensions across devices
        self.d_model_per_rank = d_model // world_size
        self.n_heads_per_rank = n_heads // world_size

        # Local attention heads
        self.local_attention = nn.MultiheadAttention(
            embed_dim=d_model,
            num_heads=self.n_heads_per_rank,
            batch_first=True
        )

        # Local feed-forward dimensions
        d_ff_per_rank = d_ff // world_size
        self.ff1 = nn.Linear(d_model, d_ff_per_rank)
        self.ff2 = nn.Linear(d_ff_per_rank, d_model)

        self.norm1 = nn.LayerNorm(d_model)
        self.norm2 = nn.LayerNorm(d_model)

    def forward(self, x: Tensor) -> Tensor:
        # Attention with all-reduce
        residual = x
        x = self.norm1(x)

        # Each rank computes its attention heads
        local_attn_out, _ = self.local_attention(x, x, x)

        # All-reduce attention outputs
        if self.world_size > 1:
            torch.distributed.all_reduce(local_attn_out)
            local_attn_out = local_attn_out / self.world_size

        x = residual + local_attn_out

        # Feed-forward with model parallelism
        residual = x
        x = self.norm2(x)

        # Split activation across ranks
        x_split = x.chunk(self.world_size, dim=-1)[self.rank]

        # Local feed-forward
        ff_out = F.gelu(self.ff1(x))
        ff_out = self.ff2(ff_out)

        # Gather feed-forward outputs
        if self.world_size > 1:
            ff_outputs = [torch.zeros_like(ff_out) for _ in range(self.world_size)]
            torch.distributed.all_gather(ff_outputs, ff_out)
            ff_out = torch.cat(ff_outputs, dim=-1)

        x = residual + ff_out
        return x
\end{lstlisting}

\subsection{Debugging and Visualization}

\subsubsection{Attention Pattern Analysis}

Tools for understanding what the model learns:

\begin{lstlisting}[language=Python, caption=Attention Visualization Tools]
import matplotlib.pyplot as plt
import seaborn as sns
from typing import Optional

class AttentionVisualizer:
    """Visualize attention patterns"""

    @staticmethod
    def plot_attention_heatmap(
        attention_weights: Tensor,
        input_tokens: list[str],
        layer_idx: int,
        head_idx: int,
        save_path: Optional[str] = None
    ) -> None:
        """Plot attention weights as heatmap"""

        # Extract specific head's attention
        attn = attention_weights[head_idx].cpu().numpy()

        # Create figure
        plt.figure(figsize=(10, 8))
        sns.heatmap(
            attn,
            xticklabels=input_tokens,
            yticklabels=input_tokens,
            cmap='Blues',
            cbar_kws={'label': 'Attention Weight'}
        )

        plt.title(f'Attention Pattern - Layer {layer_idx}, Head {head_idx}')
        plt.xlabel('Keys')
        plt.ylabel('Queries')
        plt.tight_layout()

        if save_path:
            plt.savefig(save_path, dpi=150)
        plt.show()

    @staticmethod
    def plot_attention_flow(
        attention_weights: list[Tensor],
        token_idx: int,
        input_tokens: list[str]
    ) -> None:
        """Visualize attention flow through layers for specific token"""

        n_layers = len(attention_weights)
        n_heads = attention_weights[0].shape[0]

        fig, axes = plt.subplots(n_layers, 1, figsize=(12, 3 * n_layers))

        for layer_idx, layer_attn in enumerate(attention_weights):
            # Average across heads
            avg_attn = layer_attn.mean(dim=0)[token_idx].cpu().numpy()

            ax = axes[layer_idx] if n_layers > 1 else axes
            ax.bar(range(len(input_tokens)), avg_attn)
            ax.set_xticks(range(len(input_tokens)))
            ax.set_xticklabels(input_tokens, rotation=45)
            ax.set_ylabel('Attention Weight')
            ax.set_title(f'Layer {layer_idx} - Attention from "{input_tokens[token_idx]}"')

        plt.tight_layout()
        plt.show()

    @staticmethod
    def analyze_attention_entropy(
        attention_weights: Tensor
    ) -> dict[str, Tensor]:
        """Compute attention entropy statistics"""

        # Add small epsilon to avoid log(0)
        eps = 1e-9
        attn = attention_weights + eps

        # Compute entropy
        entropy = -(attn * torch.log(attn)).sum(dim=-1)

        # Statistics
        stats = {
            'mean_entropy': entropy.mean(),
            'std_entropy': entropy.std(),
            'min_entropy': entropy.min(),
            'max_entropy': entropy.max(),
            'entropy_per_head': entropy.mean(dim=-1)
        }

        return stats

class GradientAnalyzer:
    """Analyze gradient flow in transformers"""

    def __init__(self, model: nn.Module) -> None:
        self.model = model
        self.gradient_norms: dict[str, list[float]] = {}
        self.handles = []

        # Register hooks
        for name, param in model.named_parameters():
            if param.requires_grad:
                handle = param.register_hook(
                    lambda grad, name=name: self._save_gradient(name, grad)
                )
                self.handles.append(handle)
                self.gradient_norms[name] = []

    def _save_gradient(self, name: str, grad: Tensor) -> None:
        """Save gradient norm"""
        self.gradient_norms[name].append(grad.norm().item())

    def plot_gradient_flow(self) -> None:
        """Plot gradient norms across layers"""

        # Get average gradient norms
        avg_grads = {}
        for name, grads in self.gradient_norms.items():
            if grads:
                avg_grads[name] = np.mean(grads)

        # Sort by layer order
        layers = []
        values = []
        for name in sorted(avg_grads.keys()):
            if 'weight' in name:  # Only plot weight gradients
                layers.append(name.replace('.weight', ''))
                values.append(avg_grads[name])

        # Plot
        plt.figure(figsize=(15, 5))
        plt.bar(range(len(layers)), values)
        plt.xticks(range(len(layers)), layers, rotation=90)
        plt.xlabel('Layers')
        plt.ylabel('Average Gradient Norm')
        plt.title('Gradient Flow Through Network')
        plt.yscale('log')
        plt.tight_layout()
        plt.show()

    def cleanup(self) -> None:
        """Remove hooks"""
        for handle in self.handles:
            handle.remove()
\end{lstlisting}
